\documentclass[12pt,a4paper]{article}

\usepackage[utf8]{inputenc}
\usepackage[T1]{fontenc}
\usepackage[brazil]{babel}
\usepackage{geometry}
\usepackage{booktabs}
\usepackage{array}
\usepackage{xcolor}
\usepackage{tikz}
\usepackage{float}

\usetikzlibrary{arrows.meta, positioning}

\geometry{margin=2.5cm}

\definecolor{casaClara}{HTML}{F1DFC5}
\definecolor{casaEscura}{HTML}{6D8F9C}
\definecolor{texto}{HTML}{253238}
\definecolor{destaque}{HTML}{D85C47}
\definecolor{azul}{HTML}{3379A8}

\tikzset{
    rainha/.style={
        circle,
        fill=texto,
        text=white,
        minimum size=0.72cm,
        inner sep=0pt,
        font=\bfseries
    },
    seta/.style={
        -{Latex[length=3mm]},
        line width=0.8pt,
        draw=texto
    },
    troca/.style={
        -{Latex[length=3mm]},
        line width=1.1pt,
        draw=azul
    },
    conflito/.style={
        draw=destaque,
        line width=1.4pt,
        dashed
    }
}

\newcommand{\desenharTabuleiro}{
    \foreach \linha in {0,...,3} {
        \foreach \coluna in {0,...,3} {
            \pgfmathtruncatemacro{\cor}{mod(\linha+\coluna,2)}
            \ifnum\cor=0
                \fill[casaClara] (\coluna,-\linha) rectangle ++(1,-1);
            \else
                \fill[casaEscura] (\coluna,-\linha) rectangle ++(1,-1);
            \fi
        }
    }
    \draw[texto, line width=0.7pt] (0,0) rectangle (4,-4);
    \foreach \i in {1,2,3} {
        \draw[black!20] (\i,0) -- (\i,-4);
        \draw[black!20] (0,-\i) -- (4,-\i);
    }
}

\newcommand{\rainha}[2]{
    \node[rainha] at ({#2+0.5},{-#1-0.5}) {Q};
}

\title{Simulated Annealing aplicado ao Problema das N-Rainhas}
\author{}
\date{}

\begin{document}

\maketitle

\section{Introdução}

O problema das N-Rainhas consiste em posicionar $N$ rainhas em um tabuleiro $N \times N$ de forma que nenhuma rainha ataque outra. Como rainhas atacam na mesma linha, coluna e diagonais, uma solução válida deve evitar todos esses conflitos.

Neste trabalho foi implementado o algoritmo Simulated Annealing para buscar soluções para o problema. A técnica é inspirada no processo de recozimento de metais, em que uma solução pode aceitar estados piores temporariamente para escapar de mínimos locais. Ao longo da execução, a temperatura diminui e o algoritmo passa a aceitar menos movimentos ruins.

\section{Modelagem do Problema}

A representação utilizada foi um vetor de inteiros. Cada posição do vetor representa uma linha do tabuleiro, e o valor armazenado representa a coluna onde a rainha daquela linha está posicionada. Por exemplo, em um tabuleiro $N = 4$, o vetor:

\[
[1, 3, 0, 2]
\]

representa rainhas nas posições $(0,1)$, $(1,3)$, $(2,0)$ e $(3,2)$.

A solução inicial é gerada como uma permutação aleatória dos valores de $0$ até $N-1$. Assim, cada linha possui exatamente uma rainha e cada coluna também é usada uma única vez. Isso reduz bastante o espaço de busca, pois os principais conflitos restantes ocorrem nas diagonais.

\begin{figure}[H]
\centering
\begin{tikzpicture}
    \node[font=\Large\ttfamily, text=texto] (vetor) at (0,0) {[1, 3, 0, 2]};
    \node[below=0.25cm of vetor, align=center, text=texto] {
        índice = linha\\
        valor = coluna
    };

    \draw[seta] (2.1,0) -- (4.2,0);

    \begin{scope}[shift={(4.8,1.8)}]
        \desenharTabuleiro
        \rainha{0}{1}
        \rainha{1}{3}
        \rainha{2}{0}
        \rainha{3}{2}

        \foreach \linha in {0,...,3} {
            \node[left, font=\small, text=texto] at (-0.15,{-\linha-0.5}) {\linha};
        }
        \foreach \coluna in {0,...,3} {
            \node[below, font=\small, text=texto] at ({\coluna+0.5},-4.15) {\coluna};
        }
        \node[font=\small, text=texto] at (2,-4.75) {colunas};
        \node[font=\small, rotate=90, text=texto] at (-0.75,-2) {linhas};
    \end{scope}
\end{tikzpicture}
\caption{Representação do estado por vetor.}
\end{figure}

\section{Função de Custo}

A função de custo mede a quantidade de conflitos entre rainhas. O objetivo do algoritmo é minimizar essa função até chegar ao valor zero.

No código, são verificados:

\begin{itemize}
    \item conflitos nas diagonais;
    \item conflitos de coluna, como verificação adicional.
\end{itemize}

Quanto menor o número de conflitos, melhor é a solução. Uma solução com custo igual a zero é uma solução válida para o problema.

\begin{figure}[H]
\centering
\begin{tikzpicture}
    \node[font=\Large\ttfamily, text=texto] at (0,0) {[0, 1, 3, 2]};
    \node[text=texto] at (0,-0.65) {estado com conflitos};

    \draw[seta] (1.9,0) -- (3.8,0);

    \begin{scope}[shift={(4.4,1.8)}]
        \desenharTabuleiro

        \draw[conflito] (0.5,-0.5) -- (1.5,-1.5);
        \draw[conflito] (3.5,-2.5) -- (2.5,-3.5);

        \rainha{0}{0}
        \rainha{1}{1}
        \rainha{2}{3}
        \rainha{3}{2}
    \end{scope}

    \node[
        draw=destaque,
        fill=destaque!12,
        rounded corners,
        inner sep=7pt,
        text=texto
    ] at (6.4,-3.2) {custo = 2 conflitos diagonais};
\end{tikzpicture}
\caption{Exemplo visual da função de custo contando pares de rainhas em conflito.}
\end{figure}

\section{Geração de Vizinhos}

A estratégia de vizinhança utilizada consiste em escolher duas linhas aleatórias e trocar as colunas das rainhas dessas linhas. Como o tabuleiro é representado por uma permutação, essa troca mantém uma rainha por linha e preserva o uso das colunas.

Essa vizinhança é simples, mas adequada para o problema, pois modifica parcialmente o estado atual sem reconstruir todo o tabuleiro.

\begin{figure}[H]
\centering
\begin{tikzpicture}
    \begin{scope}[font=\Large\ttfamily, text=texto]
        \node at (-1.20,0) {[};
        \node at (-0.90,0) {1};
        \node at (-0.60,0) {,};
        \node[draw=azul, rounded corners, line width=0.9pt, inner sep=2.5pt] at (-0.30,0) {3};
        \node at (0.00,0) {,};
        \node[draw=azul, rounded corners, line width=0.9pt, inner sep=2.5pt] at (0.30,0) {0};
        \node at (0.60,0) {,};
        \node at (0.90,0) {2};
        \node at (1.20,0) {]};

        \node at (5.80,0) {[};
        \node at (6.10,0) {1};
        \node at (6.40,0) {,};
        \node[draw=azul, rounded corners, line width=0.9pt, inner sep=2.5pt] at (6.70,0) {0};
        \node at (7.00,0) {,};
        \node[draw=azul, rounded corners, line width=0.9pt, inner sep=2.5pt] at (7.30,0) {3};
        \node at (7.60,0) {,};
        \node at (7.90,0) {2};
        \node at (8.20,0) {]};
    \end{scope}

    \draw[troca] (1.55,0.15) .. controls (3.2,1.0) and (4.0,1.0) .. (5.45,0.15);
    \node[font=\small, text=texto] at (3.55,1.25) {troca das linhas 1 e 2};

    \node[text=texto] at (0,-1.0) {estado atual};
    \node[text=texto] at (7,-1.0) {vizinho gerado};

    \begin{scope}[shift={(-0.6,-1.8)}, scale=0.75]
        \desenharTabuleiro
        \rainha{0}{1}
        \rainha{1}{3}
        \rainha{2}{0}
        \rainha{3}{2}
    \end{scope}

    \draw[seta] (3.1,-3.3) -- (4.35,-3.3);

    \begin{scope}[shift={(4.8,-1.8)}, scale=0.75]
        \desenharTabuleiro
        \rainha{0}{1}
        \rainha{1}{0}
        \rainha{2}{3}
        \rainha{3}{2}
    \end{scope}
\end{tikzpicture}
\caption{Geração de vizinho pela troca de duas posições do vetor.}
\end{figure}

\section{Simulated Annealing}

O algoritmo inicia com uma solução aleatória e uma temperatura inicial $T_0$. A cada iteração, é gerado um vizinho e comparado seu custo com o custo da solução atual.

Se o vizinho for melhor, ele é aceito. Se for pior, ainda pode ser aceito com probabilidade:

\[
P = e^{-\Delta E / T}
\]

onde $\Delta E$ é a diferença entre o custo do vizinho e o custo atual, e $T$ é a temperatura corrente.

Após cada iteração, a temperatura é reduzida por resfriamento geométrico:

\[
T = T \cdot \alpha
\]

Os critérios de parada utilizados foram:

\begin{itemize}
    \item encontrar uma solução com custo zero;
    \item temperatura menor que $T_{min}$;
    \item atingir o limite máximo de iterações.
\end{itemize}

Na implementação, foram usados $T_{min} = 0{,}0001$ e limite máximo de $100000$ iterações.

\section{Experimentos}

Foram realizados experimentos com os valores de $N$ sugeridos na especificação: $8$, $16$, $32$ e $128$. Para cada valor de $N$, foram testadas três combinações de temperatura inicial e taxa de resfriamento. Cada combinação foi executada três vezes, e a tabela apresenta as médias obtidas.

\begin{center}
\begin{tabular}{ccccc}
\toprule
$N$ & $T_0$ & $\alpha$ & Conflitos médios & Iterações médias \\
\midrule
8   & 20 & 0{,}9995 & 0{,}000000  & 340{,}333333 \\
8   & 5  & 0{,}9500 & 1{,}000000  & 211{,}000000 \\
8   & 1  & 0{,}9000 & 0{,}333333  & 58{,}333333 \\
\midrule
16  & 20 & 0{,}9995 & 0{,}000000  & 7081{,}666667 \\
16  & 5  & 0{,}9500 & 0{,}666667  & 210{,}000000 \\
16  & 1  & 0{,}9000 & 1{,}666667  & 88{,}000000 \\
\midrule
32  & 20 & 0{,}9995 & 0{,}000000  & 10436{,}333333 \\
32  & 5  & 0{,}9500 & 2{,}666667  & 211{,}000000 \\
32  & 1  & 0{,}9000 & 5{,}333333  & 88{,}000000 \\
\midrule
128 & 20 & 0{,}9995 & 0{,}000000  & 15422{,}000000 \\
128 & 5  & 0{,}9500 & 32{,}000000 & 211{,}000000 \\
128 & 1  & 0{,}9000 & 45{,}333333 & 88{,}000000 \\
\bottomrule
\end{tabular}
\end{center}

\begin{center}
\begin{tabular}{cccc}
\toprule
$N$ & $T_0$ & $\alpha$ & Tempo médio \\
\midrule
8   & 20 & 0{,}9995 & 159{,}224 $\mu$s \\
8   & 5  & 0{,}9500 & 89{,}444 $\mu$s \\
8   & 1  & 0{,}9000 & 27{,}618 $\mu$s \\
16  & 20 & 0{,}9995 & 4{,}305241 ms \\
16  & 5  & 0{,}9500 & 120{,}495 $\mu$s \\
16  & 1  & 0{,}9000 & 44{,}563 $\mu$s \\
32  & 20 & 0{,}9995 & 17{,}476301 ms \\
32  & 5  & 0{,}9500 & 330{,}651 $\mu$s \\
32  & 1  & 0{,}9000 & 136{,}746 $\mu$s \\
128 & 20 & 0{,}9995 & 387{,}604950 ms \\
128 & 5  & 0{,}9500 & 5{,}155863 ms \\
128 & 1  & 0{,}9000 & 2{,}108906 ms \\
\bottomrule
\end{tabular}
\end{center}

\section{Análise dos Resultados}

Os resultados mostram que a combinação $T_0 = 20$ e $\alpha = 0{,}9995$ foi a mais consistente. Ela encontrou soluções com zero conflitos para todos os valores de $N$ testados. Isso ocorreu porque a temperatura inicial maior e o resfriamento mais lento permitiram ao algoritmo explorar melhor o espaço de busca antes de se tornar mais restritivo.

As combinações com $T_0 = 5$, $\alpha = 0{,}95$ e $T_0 = 1$, $\alpha = 0{,}9$ foram mais rápidas, mas apresentaram mais conflitos, principalmente para $N = 32$ e $N = 128$. Isso indica que o resfriamento mais rápido reduz o tempo de execução, porém aumenta a chance de o algoritmo parar em uma solução ainda não ideal.

Também é possível perceber que, conforme $N$ aumenta, o tempo de execução cresce. Isso acontece porque a função de custo compara pares de rainhas para detectar conflitos, o que torna cada avaliação mais cara em tabuleiros maiores.

\section{Conclusão}

A implementação conseguiu aplicar o Simulated Annealing ao problema das N-Rainhas e produzir soluções válidas nos principais testes. O experimento mostrou a importância da escolha dos parâmetros: temperaturas maiores e resfriamento mais lento tendem a melhorar a qualidade da solução, enquanto temperaturas menores e resfriamento rápido reduzem o tempo, mas podem deixar conflitos no tabuleiro.

Assim, o trabalho confirma que o Simulated Annealing é uma abordagem adequada para problemas de busca heurística, especialmente quando é necessário escapar de mínimos locais e explorar um espaço de soluções muito grande.

\end{document}
